{{دیوان
| عنوان = آزادی رائے
| مصنف = سید احمد خان
| صنف = مضمون
| ماخذ = ویکی ماخذ
| ماخذ_ربط = https://ur.wikisource.org/wiki/%D8%A2%D8%B2%D8%A7%D8%AF%DB%8C_%D8%B1%D8%A7%D8%A6%DB%92
}}
== آزادی رائے ==

ہم اپنے اس آرٹیکل کو ایک بڑے لائق اور قابل زمانہ حال کے [[فیلسوف]] کی تحریر سے اخذ کرتے ہیں۔ رائے کی آزادی ایک ایسی چیز ہے کہ ہر ایک انسان اس پر پورا پورا حق رکھتا ہے۔ فرض کرو کہ تمام آدمی بجز ایک شخص کے کسی بات پر متفق الرائے ہیں، مگر صرف وہی ایک شخص ان کے بر خلاف رائے رکھتا ہے تو ان تمام آدمیوں کو اس ایک شخص کی رائے کو غلط ٹھہرانے کے لیے اس سے زیادہ کچھ استحقاق نہیں ہے جتنا کہ اس ایک شخص کو ان تمام آدمیوں کی رائے کے غلط ثابت کرنے کا (اگر وہ ثابت کر سکے) استحقاق حاصل ہے۔<ref>مثال کے لیے حاشیہ: یہ مضمون تہذیب الاخلاق میں شائع ہوا۔</ref>

کوئی وجہ اس بات کی نہیں ہے کہ پانچ آدمیوں کو تو بمقابلہ پانچ آدمیوں کی رایوں کے غلط ٹھہرانے کا استحقاق ہو اور ایک آدمی کو بمقابل نو آدمیوں کے یہ استحقاق نہ ہو۔ رائے کی غلطی آدمیوں کی تعداد کی کمی بیشی پر منحصر نہیں ہے، بلکہ قوت استدلال پر منحصر ہے۔ جیسے کہ یہ بات ممکن ہے کہ نو آدمیوں کی رائے بمقابلہ ایک شخص کے صحیح ہو ویسے یہ بھی ممکن ہے کہ ایک شخص کی رائے بمقابل نو کے صحیح ہو۔

رایوں کا بند رہنا خواہ بسبب کسی مذہبی خوف کے اور خواہ بسبب اندیشہ برادری وقوم کے اور خواہ بدنامی کے ڈر سے اور یا گورنمنٹ کے ظلم سے نہایت ہی بری چیز ہے۔ اگر رائے اس قسم کی کوئی چیز ہوتی جس کی قدر و قیمت صرف اس رائے والے کی ذات ہی سے متعلق اور اسی میں محصور ہوتی تو رایوں کے بند رہنے سے ایک خاص شخص یا معدودے چند کا نقصان متصور ہوتا مگر رایوں کے بند رہنے سے تمام انسانوں کی حق تلفی ہوتی ہے اور کل انسانوں کو نقصان پہنچتا ہے اور نہ صرف موجودہ انسانوں کو، بلکہ ان کو بھی جو آئندہ پیدا ہوں گے۔
